\chapter*{List of Symbols and Abbreviations}
\addcontentsline{toc}{chapter}{List of Symbols and Abbreviations}

\begin{longtable}{@{}p{0.2\textwidth}p{0.74\textwidth}@{}}
\toprule
Symbol & Meaning\\
\midrule
\endhead
$\F,\ \Fbar$ & microscopic and macroscopic deformation gradient\\
$J=\det\F$ & volume ratio\\
$\Cten=\F^{\mathsf T}\F$ & right Cauchy--Green tensor\\
$I_1,\ I_2$ & first and second principal invariants of $\Cten$\\
$\Psi,\ \Psibar$ & microscopic and homogenized strain-energy density\\
$\Pstress,\ \Pbar$ & microscopic and homogenized first Piola--Kirchhoff stress\\
$\Aten,\ \Abar$ & microscopic and homogenized material tangent $\partial\Pstress/\partial\F$ (stored as $9\times9$)\\
$\bm\sigma,\ p,\ \bm s$ & Cauchy stress, pressure $\tfrac13\operatorname{tr}\bm\sigma$, deviator $\bm\sigma-p\Iten$\\
$\bm w$ & periodic displacement fluctuation in the RVE\\
$\Rten\in SO(3)$ & particle rotation\\
$\vfrac$ & particle volume fraction\\
$\mu,\ \kappa$ & matrix shear and bulk modulus ($\mu=1$, $\kappa=10^4$)\\
$\mu_{\mathrm{eff}},\ K_{\mathrm{eff}}$ & effective small-strain shear and bulk modulus\\
$h$ & maximum element size of the macroscopic mesh\\
$\bm X,\ \bm x=\bm\varphi(\bm X)$ & reference and current position; $\bm\varphi$ is the deformation map\\
$\Omega_0,\ \Omega$ & reference body of the macroscopic problem; RVE (unit cube)\\
$\bm u,\ \delta\bm u$ & displacement and test function\\
$\bar I_1,\ \bar I_2$ & invariants of the isochoric part: $I_1=J^{2/3}\bar I_1$, $I_2=J^{4/3}\bar I_2$\\
$\bm\varepsilon$ & small-strain tensor\\
$\Pi,\ p$ & mixed (displacement--pressure) functional; Lagrange pressure\\
$\lambda$ & Lam\'e constant of the neo-Hookean reference law ($\lambda=10$, \cref{ch:theory}); in \cref{ch:methods} also an eigenvalue ($\lambda_{\min}(\bm K)$)\\
$K_{\mathrm{MT}},\ \mu_{\mathrm{MT}},\ \eta$ & Mori--Tanaka bulk and shear modulus; auxiliary modulus in \cref{eq:mt}\\
$\bm y,\ \hat{\bm y},\ \hat c$ & invariant features $[I_1,I_2,J,-\ln J]$, their standardized values, standardized volume fraction\\
$W^{+}=\mathrm{softplus}(W)$ & non-negative network weights (the symbol $\sigma$ also denotes the softplus activation in the network equations)\\
$\bm\beta,\ \bm b_l$ & biases of the energy network; bias shifts produced by the context network\\
$\Psi_{\mathrm{raw}},\ \widehat\Psi$ & network output and its normalized value (zero energy and stress at $\F=\Iten$)\\
$\mathcal L,\ w_P,\ s_\Psi,\ s_P$ & training loss, stress weight, energy and stress scales\\
$\Psi_{\mathrm{iso}},\ \Psi_{\mathrm{vol}}$ & isochoric and volumetric energy of the proposed split\\
$\bm K,\ \bm R,\ \alpha$ & tangent stiffness matrix, residual, spectral shift of the Newton fallback\\
$\bm\theta$ & rotation increment of the particle\\
$\gamma$ & amount of simple shear\\
$\epsilon,\ \varepsilon_{\mathrm{mach}}$ & finite-difference step; machine epsilon\\
$\rho$ & Spearman rank correlation\\
\midrule
Abbreviation & Meaning\\
\midrule
AD & automatic differentiation\\
DOF & degree of freedom\\
FD & finite difference\\
MT & Mori--Tanaka\\
NH, SVK & neo-Hookean, Saint~Venant--Kirchhoff (reference material laws)\\
OCC & Open CASCADE geometry kernel used by Netgen\\
PBS & portable batch system (cluster job scheduler)\\
RMS & root mean square\\
$SO(3)$ & group of three-dimensional rotations\\
FE, FEM & finite element(s), finite-element method\\
FE$^2$ & concurrent two-scale finite-element method\\
ICNN / pICNN & (partially) input-convex neural network\\
NRMSE & normalized root-mean-square error\\
PBC & periodic boundary condition\\
PDMS & polydimethylsiloxane\\
RVE & representative volume element\\
$P_1,\ P_2$ & linear / quadratic Lagrange elements\\
\bottomrule
\end{longtable}
